% SPDX-License-Identifier: Apache-2.0
\section{One Ethernet Wire, Two Users}
\label{sec:x-ethshare}

The board has one Ethernet port and two things that want it: the
remote peripheral, whose transport is one frame per transaction, and
the port a Zephyr driver talks to through \code{EthSlots}
(\S\ref{sec:x-ethslots}). Neither can have the wire to itself.

Nothing new had to be invented to share it. Ethernet already carries
several protocols on one wire and says which is which: bytes twelve
and thirteen of every frame are its EtherType. The remote peripheral
sends \code{0x88b5}, which IEEE leaves for local use, and Zephyr's
stack sends only the standard types, so the two cannot be confused.
\code{EthShare} splits arriving frames by that type and merges
departing ones a frame at a time.

\subsection{Why the header is held}

A frame's user is not known until its fourteenth byte, and the
thirteen before it belong to that user too. So the unit holds the
first fourteen bytes, reads the type from the last two of them, and
then hands the whole frame on, held bytes first.

\lstinputlisting[rangebeginprefix=//\ begin\{,rangebeginsuffix=\},%
rangeendprefix=//\ end\{,rangeendsuffix=\},includerangemarker=false,%
linerange=state-state,caption={\code{lib/parts/src/ethshare.rs}: what
the unit holds while it learns whose a frame is.},%
label={lst:x-ethshare}]{lib/parts/src/ethshare.rs}

\subsection{The three edges, and a test that passed by luck}

Every assertion passed on the first run, which is when a test is
worth breaking on purpose. Three cases were, each with the unit
changed to be wrong about it and nothing else.

\emph{A frame of exactly fourteen bytes} ends on the second byte of
its own type, so the unit learns its user and that it is over in the
same cycle. Marking no last byte there made the unit wait for the
rest of a frame that had ended; that frame never arrived, and the run
failed with it missing.

\emph{Two users sending at once} must take turns. Letting the first
always go first sent both its frames before either of the second's,
and the run failed on the order.

\emph{A frame that ends before its type is known} is the one that
taught something. The unit must not read a type it has not reached,
and the first version of the run passed with that check removed. The
type it would have read was the high byte left behind by the frame
before, \code{0x88}, beside the short frame's last byte; that byte
happened to be \code{0xc9}, so the stale type was \code{0x88c9}, it
matched nothing, and the frame went the right way by luck of its
data.

So the short frame is now built so that its last byte is
\code{0xb5}. A unit reading a type it has not reached sees
\code{0x88b5}, sends the frame to the wrong user, and the run fails:
the first user takes four frames and the second one. A test is only
adversarial when its data is chosen to be, and this one was not until
it was broken on purpose and did not notice.

\begin{center}\footnotesize
\begin{tabular}{@{}lll@{}}
\toprule
The unit made wrong about & first run & now \\
\midrule
a frame ending on its type & fails & fails \\
taking turns & fails & fails \\
a type it has not reached & \textbf{passes} & fails \\
\bottomrule
\end{tabular}
\end{center}

\lstinputlisting[style=ascii,caption={What \code{ex\_ethshare}
printed when the build ran it.},%
label={lst:x-ethshare-out}]{out_ethshare.txt}
